\section{Self-dual channels}
\label{sec:self_dual_channels}
%% =========================================================

The transfer matrix $\widehat T_{\alpha\beta}=\tr(F_\alpha^\dagger T(G_\beta))$
of \cite[Section~2.2, equation~(2.20)]{Wolf2012Quantum} is built from two
orthonormal families. Taking one and the same family on both sides turns the
Hermiticity of $\widehat T$ into a statement about $T$ alone, and that statement
is self-duality. Orthonormality is always meant for the Hilbert--Schmidt
pairing, $\tr(G_\beta^\dagger G_{\beta'})=\delta_{\beta\beta'}$
\cite[Section~2.2]{Wolf2012Quantum}; the families are not assumed Hermitian.

\begin{definition}[Hilbert--Schmidt orthonormal family]
    \label{def:hilbert_schmidt_orthonormal}
    \lean{HilbertSchmidtOrthonormal}
    \leanok
    A family $(\sigma_\alpha)$ in $\MN{d}$ is \emph{Hilbert--Schmidt orthonormal}
    when
    \begin{align}
        \tr\bigl(\sigma_\alpha^\dagger \sigma_\beta\bigr)
        &= \delta_{\alpha\beta}.
        \notag
    \end{align}
\end{definition}

\begin{definition}[Transfer matrix in one orthonormal family]
    \label{def:transfer_matrix_of_basis}
    \lean{transferMatrixOfBasis}
    \leanok
    For a linear map $T:\MN{d}\to\MN{d}$ and a family $(\sigma_\alpha)$ in
    $\MN{d}$ used on both sides, the \emph{transfer matrix of $T$ in
    $(\sigma_\alpha)$} is
    \begin{align}
        \widehat T_{\alpha\beta}
        &= \tr\bigl(\sigma_\alpha^\dagger T(\sigma_\beta)\bigr).
        \notag
    \end{align}
\end{definition}

\begin{lemma}[Coordinates in a Hilbert--Schmidt orthonormal basis]
    \label{lem:hilbert_schmidt_orthonormal_coordinates}
    \lean{HilbertSchmidtOrthonormal.repr_apply,
        HilbertSchmidtOrthonormal.eq_zero_of_trace_conjTranspose_mul,
        HilbertSchmidtOrthonormal.tracePairingSelfDualBasis}
    \leanok
    \uses{def:hilbert_schmidt_orthonormal}
    Let $(\sigma_\alpha)$ be a Hilbert--Schmidt orthonormal basis of $\MN{d}$.
    The coordinates of $X\in\MN{d}$ in this basis are
    $\tr(\sigma_\alpha^\dagger X)$; consequently $X=0$ as soon as
    $\tr(\sigma_\alpha^\dagger X)=0$ for every $\alpha$. If in addition every
    $\sigma_\alpha$ is Hermitian, the coordinates are $\tr(\sigma_\alpha X)$, so
    the basis is self-dual for the bilinear trace pairing.
\end{lemma}

\begin{proof}\leanok
    Expanding $X=\sum_\beta c_\beta\sigma_\beta$ and pairing with
    $\sigma_\alpha$ gives $\tr(\sigma_\alpha^\dagger X)=c_\alpha$ by
    orthonormality. Vanishing of all coordinates forces $X=0$. For Hermitian
    $\sigma_\alpha$ the two pairings $\tr(\sigma_\alpha^\dagger X)$ and
    $\tr(\sigma_\alpha X)$ agree.
\end{proof}

\begin{lemma}[Transfer matrix of the dual map]
    \label{lem:transfer_matrix_of_basis_dual}
    \lean{transferMatrixOfBasis_conjTranspose,
        transferMatrix_conjTranspose_eq_traceAdjointMap}
    \leanok
    \uses{def:transfer_matrix_of_basis, def:trace_pairing_adjoint,
        def:mpu_transfer_matrix}
    Let $T:\MN{d}\to\MN{d}$ satisfy $T(X^\dagger)=T(X)^\dagger$. Then in any
    family $(\sigma_\alpha)$ used on both sides,
    \begin{align}
        \bigl(\widehat T\bigr)^\dagger &= \widehat{T^*},
        \notag
    \end{align}
    and the same identity holds for the transfer matrix taken in the matrix
    units.
\end{lemma}

\begin{proof}\leanok
    \uses{def:trace_pairing_adjoint, def:mpu_transfer_matrix}
    The $(\alpha,\beta)$ entry of $(\widehat T)^\dagger$ is
    $\overline{\tr(\sigma_\beta^\dagger T(\sigma_\alpha))}
    =\tr\bigl(T(\sigma_\alpha)^\dagger\sigma_\beta\bigr)$. Hermiticity
    preservation replaces $T(\sigma_\alpha)^\dagger$ by
    $T(\sigma_\alpha^\dagger)$, and the defining identity
    $\tr(AT^*(B))=\tr(T(A)B)$ of the dual map turns the result into
    $\tr(\sigma_\alpha^\dagger T^*(\sigma_\beta))$.

    In the matrix units the same two steps appear entrywise. There the
    $\bigl((j,i),(\ell,k)\bigr)$ entry of $\widehat T$ is
    $\bigl(T(E_{k\ell})\bigr)_{ij}$, and the entry of
    $\widehat{T^*}$ is $\tr\bigl(E_{k\ell}\,T(E_{ji})\bigr)$ by the entrywise
    form of the dual map; conjugating the first and using
    $E_{k\ell}^\dagger=E_{\ell k}$ together with Hermiticity preservation
    identifies the two.
\end{proof}

\begin{lemma}[Dual map of a Kraus family]
    \label{lem:trace_adjoint_of_kraus}
    \lean{traceAdjointMap_of_kraus, map_conjTranspose_of_kraus,
        IsCPMap.map_conjTranspose}
    \leanok
    \uses{def:trace_pairing_adjoint, def:cp_map}
    If $T(X)=\sum_j K_j X K_j^\dagger$ then $T^*(X)=\sum_j K_j^\dagger X K_j$: the
    Kraus operators of $T$ and of $T^*$ differ by Hermitian conjugation. Such a
    $T$ also satisfies $T(X^\dagger)=T(X)^\dagger$.
\end{lemma}

\begin{proof}\leanok
    For every $N$,
    $\tr\bigl(\sum_j K_j^\dagger XK_j\,N\bigr)
    =\sum_j\tr\bigl(X\,K_jNK_j^\dagger\bigr)=\tr(X\,T(N))$ by cyclicity of the
    trace, which is the defining identity of $T^*$. Nondegeneracy of the trace
    pairing identifies the two maps. Hermiticity preservation is the
    conjugate transpose of the Kraus sum.
\end{proof}

\begin{lemma}[Hermitian recombination of a symmetrized Kraus family]
    \label{lem:hermitian_kraus_of_self_dual}
    \lean{exists_hermitian_kraus_of_self_dual}
    \leanok
    Suppose $T(X)=\sum_j K_j X K_j^\dagger$ and also
    $T(X)=\sum_j K_j^\dagger X K_j$. Then $T$ has a family of Hermitian Kraus
    operators, namely $\tfrac12(K_j+K_j^\dagger)$ and
    $\tfrac{i}{2}(K_j-K_j^\dagger)$.
\end{lemma}

\begin{proof}\leanok
    \uses{thm:kraus_same_map_of_unitary_combination}
    Averaging the two representations gives
    \begin{align}
        T(X)
        &= \frac12\sum_j\bigl(K_jXK_j^\dagger+K_j^\dagger XK_j\bigr),
        \notag
    \end{align}
    that is, the Kraus family $(K_j,K_j^\dagger)/\sqrt2$. Mixing each such pair by
    the unitary $(a,b)\mapsto\bigl((a+b)/\sqrt2,\;i(a-b)/\sqrt2\bigr)$ leaves the
    map unchanged and produces the pair $\tfrac12(K_j+K_j^\dagger)$,
    $\tfrac{i}{2}(K_j-K_j^\dagger)$, both of which are Hermitian.
\end{proof}

\begin{lemma}[The transfer matrix determines the map]
    \label{lem:transfer_matrix_of_basis_injective}
    \lean{transferMatrixOfBasis_injective}
    \leanok
    \uses{def:transfer_matrix_of_basis, def:hilbert_schmidt_orthonormal}
    Let $(\sigma_\alpha)$ be a Hilbert--Schmidt orthonormal basis of $\MN{d}$.
    For linear maps $S,T:\MN{d}\to\MN{d}$,
    \begin{align}
        \widehat T=\widehat S
        &\implies T=S.
        \notag
    \end{align}
\end{lemma}

\begin{proof}\leanok
    \uses{lem:hilbert_schmidt_orthonormal_coordinates}
    Equality of the transfer matrices reads
    $\tr\bigl(\sigma_\alpha^\dagger(T-S)(\sigma_\beta)\bigr)=0$ for all
    $\alpha,\beta$. For fixed $\beta$ this makes every coordinate of
    $(T-S)(\sigma_\beta)$ vanish, so $(T-S)(\sigma_\beta)=0$; a linear map
    vanishing on a basis is zero.
\end{proof}

\begin{lemma}[Self-duality against Hermiticity of the transfer matrix]
    \label{lem:transfer_matrix_of_basis_hermitian_iff}
    \lean{transferMatrixOfBasis_isHermitian_iff}
    \leanok
    \uses{def:transfer_matrix_of_basis, def:hilbert_schmidt_orthonormal,
        def:trace_pairing_adjoint}
    Let $T:\MN{d}\to\MN{d}$ satisfy $T(X^\dagger)=T(X)^\dagger$ and let
    $(\sigma_\alpha)$ be a Hilbert--Schmidt orthonormal basis of $\MN{d}$ used on
    both sides. Then
    \begin{align}
        \widehat T=\widehat T^\dagger
        &\iff T=T^*.
        \notag
    \end{align}
\end{lemma}

\begin{proof}\leanok
    \uses{lem:transfer_matrix_of_basis_dual,
        lem:transfer_matrix_of_basis_injective}
    By Lemma~\ref{lem:transfer_matrix_of_basis_dual},
    $\widehat T^\dagger=\widehat{T^*}$. So
    $\widehat T=\widehat T^\dagger$ reads $\widehat{T^*}=\widehat T$, which by
    Lemma~\ref{lem:transfer_matrix_of_basis_injective} is $T^*=T$; the converse
    substitutes $T^*=T$ into the same identity.
\end{proof}

\begin{lemma}[The same equivalence in the matrix units]
    \label{lem:transfer_matrix_units_hermitian_iff}
    \lean{transferMatrix_isHermitian_iff_traceAdjointMap_eq_self}
    \leanok
    \uses{def:cp_map, def:trace_pairing_adjoint, def:mpu_transfer_matrix}
    Let $T:\MN{d}\to\MN{d}$ be completely positive and let $\widehat T$ be its
    transfer matrix in the matrix units. Then
    \begin{align}
        \widehat T=\widehat T^\dagger
        &\iff T=T^*.
        \notag
    \end{align}
\end{lemma}

\begin{proof}\leanok
    \uses{lem:transfer_matrix_of_basis_dual, lem:trace_adjoint_of_kraus}
    Complete positivity gives $T(X^\dagger)=T(X)^\dagger$, so
    $\widehat T^\dagger=\widehat{T^*}$ for the matrix-unit transfer matrix as
    well, and that transfer matrix determines the map.
\end{proof}

\begin{proposition}[Self-dual channels]
    \label{prop:self_dual_channels}
    \lean{self_dual_tfae}
    \leanok
    \uses{def:cp_map, def:trace_pairing_adjoint, def:transfer_matrix_of_basis,
        def:hilbert_schmidt_orthonormal}
    Let $T:\MN{d}\to\MN{d}$ be completely positive and let $(\sigma_\alpha)$ be a
    Hilbert--Schmidt orthonormal basis of $\MN{d}$. The following are equivalent.
    \begin{enumerate}
        \item $T=T^*$.
        \item $\widehat T=\widehat T^\dagger$, the transfer matrix being taken in
            the family $(\sigma_\alpha)$ on both sides.
        \item $T$ admits a family of Hermitian Kraus operators.
    \end{enumerate}
    This is the proposition ``Self-dual channels'' of
    \cite[Section~2.2]{Wolf2012Quantum}.
\end{proposition}

\begin{proof}\leanok
    \uses{lem:transfer_matrix_of_basis_hermitian_iff, lem:trace_adjoint_of_kraus,
        lem:hermitian_kraus_of_self_dual}
    Conditions~1 and~2 are equivalent by
    Lemma~\ref{lem:transfer_matrix_of_basis_hermitian_iff}, complete positivity
    supplying the Hermiticity preservation it requires. For $3\Rightarrow1$, a
    family of Hermitian Kraus operators is fixed by the exchange
    $K_j\mapsto K_j^\dagger$ that carries $T$ to $T^*$:
    \begin{align}
        T^*(X)
        &= \sum_j K_j^\dagger XK_j
        = \sum_j K_jXK_j^\dagger
        = T(X).
        \notag
    \end{align}
    For $1\Rightarrow3$, condition~1 turns any Kraus representation
    $T(X)=\sum_jK_jXK_j^\dagger$ into a second one,
    \begin{align}
        T(X)
        &= T^*(X)
        = \sum_j K_j^\dagger XK_j,
        \notag
    \end{align}
    and the Hermitian recombination of the symmetrized family
    (Lemma~\ref{lem:hermitian_kraus_of_self_dual}) yields condition~3.
\end{proof}
